\section{Entire all-field node calculation: regularization reverses one lift inequality}
\textit{Supports: G331, n1, n3.}

Over ANY field $k$, take
\[
R=k[x,y,z]_{(x,y,z)}/(xy),\quad S=A_1\times A_2,
\quad A_1=k[x,z]_{(x,z)},\quad A_2=k[y,z]_{(y,z)}.
\]
The normalizations are the two regular planes. With $T=k[z]_{(z)}$, restriction at $x=0,y=0$ gives $R=A_1\times_T A_2$. Local denominators can be lifted with the same nonzero constant term, proving the localized fibre product as well. $S$ is generated as an $R$-module by its two idempotents, and the conductor is $\mathfrak c=(x,y)$: a conductor element must have zero common residue in $T$, and that condition suffices. Hence $R/\mathfrak c=T$, $S/\mathfrak c=T\times T$. There is exactly one nontrivial full boundary fibre: $w_1$ killing $x$ on $A_1$ and $w_2$ killing $y$ on $A_2$, both restricting to $\operatorname{ord}_z$ on $T$. Call its base point $v_0$. All other fibres are singletons. Gauge is $\mathfrak m$.

The bounded finite continuous obstacle is
\[
f(v)=-1-1/v(y),\qquad1/\infty=0.
\]
Normalized values have $v(y)\ge1$. On $A_1$ the obstacle is $-1$, and on $A_2$ it lies in $[-2,-1]$. If $o_i$ is ordinary regular order on plane $i$, every normalized valuation $w$ has $w(J)\ge o_i(J)$ for primary ideals: $J\subset\mathfrak m_i^{o_i(J)}$ and $w(\mathfrak m_i)=1$. This gives $u(w)\le u(o_i)$ for every $u\in C_{A_i}$, and then for every actual $D$ limit, including bottom. At $o_2$, $f=-2$. Thus the unconstrained source envelope is the ACTUAL $C_S$ function
\[
p_S(f\circ\rho)=E_S(f\circ\rho)=(-1,-2).              \tag{13}
\]
The constants are semilocal primary profiles (use $\mathfrak m_1\times\mathfrak m_2^2$), and are below the obstacle. The upper bounds just proved show they are greatest $D$ competitors.

For each integer $N\ge2$ put $\alpha=1/(N-1)$ and
\[
h_{1,N}(w)=-(1-\alpha)-\alpha\min\{w(x),Nw(z)\}.
\]
This is a $C_{A_1}$ profile using the proper PRIMARY ideal $(x,z^N)$, since $\min(w(x),w(z))=1$ implies the displayed minimum is at least one. It is at most $-1$, and at $w_1$ it equals $-2$. Pair it with constant $-2$ on $A_2$. The tuple agrees on the ENTIRE shared fibre, so G331 descends it to a $C_R$ competitor for $f$. At every first-plane point off $w_1$, $w(x)<\infty$; otherwise the quotient is the one-dimensional regular $T$, whose normalized section is the singleton $w_1$. At such a point the displayed profile tends to $-1$ as $N\to\infty$. On the whole second plane every downstairs $D$ competitor is at most $-2$ by regular order domination. Therefore, pointwise and using both raw $C$ and $D$ suprema,
\[
p_R f=-1\text{ on the first plane off }v_0,
\qquad p_R f=-2\text{ on the whole second plane, including }v_0. \tag{14}
\]

Monomial valuations $w_M(x)=M,w_M(z)=1$ converge to $w_1$ on EVERY polynomial and local denominator: transverse terms gain unbounded value, the surviving $z$ polynomial has its ordinary order, and denominators with nonzero constant have value zero. This works in every characteristic. First-plane values in (14) therefore raise the regularization at $v_0$ to $-1$. Away from $v_0$ on the second plane the image of the compact first plane is closed and absent, so a neighborhood avoids it. Consequently
\[
E_R f=-1\text{ on the entire first plane including }v_0,
\qquad E_R f=-2\text{ on the second plane off }v_0.    \tag{15}
\]
Its pullback is USC on each plane (a raised value at the closed point $w_2$ on the second), hence the function is USC by the compact quotient. But it is NOT in $D_R$: its restriction to $A_2$ has value $-1$ at $w_2$ and $-2$ at $o_2$, violating the proved domination inequality. At the second lift,
\[
\rho^*E_R f(w_2)=-1>-2=E_S(f\circ\rho)(w_2).          \tag{16}
\]
This is a genuine gap ONE for the SPECIFIED regularized operator, not merely the raw envelope.

The constrained upstairs raw supremum is the pullback of (14). Its ordinary regularization is $(-1,-2)$: on the first plane $w_1$ is raised to $-1$, whereas the second is constantly $-2$. Thus $E_{S,\mathrm{fib}}$ is positive and continuous but NOT fibre constant. Its full fibre maximum is exactly (15), verifying (11) at every point and the necessity of BOTH positivity and fibre constancy. On the generic branches the degree-one trace of (13) has incompatible limits $-1$ and $-2$ at the shared point, so an unrestricted continuous degree trace cannot be fabricated. In base change,
$S\otimes_R S=A_1\times T\times T\times A_2$;
the two vertical $T$ components are real parts of the support/fibre geometry and must remain.
